\documentclass[conference]{IEEEtran}

\usepackage[utf8]{inputenc}
\usepackage[T1]{fontenc}
\usepackage{amsmath,amssymb}
\usepackage{graphicx}
\usepackage{booktabs}
\usepackage{hyperref}
\usepackage{cite}

\title{RideSmart: Modelagem e An\'alise de Rotas Urbanas com Grafos: Compara\c{c}\~ao de Algoritmos de Caminho M\'{\i}nimo}

\author{
\IEEEauthorblockN{Eugenio Vitor Lopes dos Santos}
\IEEEauthorblockA{Universidade Federal do Rio Grande do Norte\\ eugenio.santos@ufrn.edu.br}
\and
\IEEEauthorblockN{Lucas Augusto da Silva Cardoso}
\IEEEauthorblockA{Universidade Federal do Rio Grande do Norte\\ lucas.cardoso@ufrn.edu.br}
\and
\IEEEauthorblockN{Pedro Henrique Ribeiro de Lima}
\IEEEauthorblockA{Universidade Federal do Rio Grande do Norte\\ pedro.lima@ufrn.edu.br}
}

\begin{document}

\maketitle

\begin{abstract}
O roteamento urbano requer a sele\c{c}\~ao otimizada de pontos de embarque que equilibrem dist\^ancia percorrida a p\'e e tempo de deslocamento de carro. Este trabalho apresenta a modelagem de uma malha vi\'aria real como grafo dirigido, obtida via OSMnx a partir do OpenStreetMap, e compara quatro algoritmos de caminho m\'{\i}nimo, a saber: Dijkstra simples, Dijkstra com fila de prioridade, A* com heur\'{\i}stica geogr\'afica e Bellman-Ford, sob diferentes condi\c{c}\~oes de tr\^ansito. Os resultados experimentais indicam que [DADOS, como algoritmo mais eficiente, tempo m\'edio e ganho com caminhada].
\end{abstract}

\begin{IEEEkeywords}
grafos, caminho m\'{\i}nimo, Dijkstra, A*, Bellman-Ford, OSMnx, mobilidade urbana
\end{IEEEkeywords}

%% ============================================================
%% I. INTRODUÇÃO
%% ============================================================
\section{Introdu\c{c}\~ao}

A mobilidade urbana enfrenta desafios crescentes de congestionamento e inefici\^encia no uso da infraestrutura vi\'aria. Aplicativos de mobilidade precisam resolver, em tempo real, o problema de roteamento que equilibra dist\^ancia percorrida a p\'e com tempo de deslocamento motorizado.

Seja $A$ o ponto de origem, $B$ o destino e $X$ a dist\^ancia m\'axima que o usu\'ario aceita caminhar. O objetivo \'e selecionar um ponto de embarque $P$, tal que $d(A, P) \le X$, que minimize o custo total da viagem composta por dois trechos:

\begin{equation}
C(P) = t_{\text{walk}}(A \to P) + t_{\text{car}}(P \to B)
\label{eq:custo}
\end{equation}

onde $t_{\text{walk}}(A \to P) = \frac{d_{\text{length}}(A \to P)}{v_{\text{walk}}}$ \'e o tempo de caminhada derivado da dist\^ancia real ao longo das ruas.

Este trabalho compara quatro algoritmos de caminho m\'{\i}nimo aplicados \`a malha vi\'aria real de [CIDADE], avaliando desempenho computacional e qualidade das solu\c{c}\~oes sob tr\^ansito sint\'etico.

%% ============================================================
%% II. MODELAGEM DO PROBLEMA
%% ============================================================
\section{Modelagem do Problema}

\subsection{Representa\c{c}\~ao do Grafo}

O grafo dirigido com multiarestas $G = (V, E)$ \'e obtido via OSMnx a partir do OpenStreetMap para a rede \texttt{drive} (ve\'iculos) e \texttt{walk} (pedestres):

\begin{itemize}
    \item \textbf{N\'os ($V$):} cruzamentos e extremidades de vias, cada qual com coordenadas geogr\'aficas (latitude, longitude).
    \item \textbf{Arestas ($E$):} segmentos de rua entre dois cruzamentos. O tipo \texttt{MultiDiGraph} admite arestas paralelas (duas ruas distintas ligando o mesmo par de n\'os) e modela ruas de m\~ao \u\'nica.
\end{itemize}

\subsection{Fun\c{c}\~oes de Custo}

Cada aresta possui at\'e tres pesos distintos, conforme apresentado na TABELA~I.

\begin{table}[htbp]
\centering
\caption{Pesos das Arestas}
\label{tab:pesos}
\begin{tabular}{@{}lll@{}}
\toprule
\textbf{Peso} & \textbf{Significado} & \textbf{Uso} \\ \midrule
\texttt{length} & dist\^ancia (m) & menor dist\^ancia e caminhada $A \to P$ \\
\texttt{travel\_time} & tempo (s) & rota mais r\'apida \\
\texttt{travel\_time\_traffic} & tempo com congestionamento & rota com tr\^ansito \\ \bottomrule
\end{tabular}
\end{table}

\subsection{Restri\c{c}\~oes e Deriva\c{c}\~oes}

O tempo de caminhada \'e calculado como $t_{\text{walk}}(A \to P) = d_{\text{length}}(A \to P) / v_{\text{walk}}$, onde $v_{\text{walk}} \approx 1{,}4$ m/s ($\approx 5$ km/h). O custo total \'e dado pela Equa\c{c}\~ao~\ref{eq:custo}, sujeito \`a restri\c{c}\~ao $d_{\text{length}}(A \to P) \le X$.

%% ============================================================
%% III. ALGORITMOS UTILIZADOS
%% ============================================================
\section{Algoritmos Utilizados}

\subsection{Dijkstra Simples}

O algoritmo de Dijkstra \cite{dijkstra} encontra o caminho m\'{\i}nimo em grafos com pesos n\~ao negativos. A implementa\c{c}\~ao via NetworkX (\texttt{nx.dijkstra\_path}) utiliza heap bin\'ario com complexidade $O((V + E) \log V)$.

\subsection{Dijkstra com Fila de Prioridade}

Implementa\c{c}\~ao manual utilizando a fila de prioridade \texttt{heapq} do Python. Cada n\'o \'e expandido em ordem de menor dist\^ancia conhecida, resultando em complexidade $O((V + E) \log V)$.

\subsection{A* com Heur\'{\i}stica Geogr\'afica}

O algoritmo A* \cite{hart} estende Dijkstra com uma heur\'{\i}stica admiss\'{\i}vel e consistente:

\begin{equation}
h(n) = \| \mathbf{pos}(n) - \mathbf{pos}(B) \|_2
\label{eq:heuristica}
\end{equation}

onde $\| \cdot \|_2$ denota a dist\^ancia euclidiana entre o n\'o $n$ e o destino $B$. A complexidade permanece $O((V + E) \log V)$, mas o n\'umero de n\'os expandidos \'e tipicamente menor que no Dijkstra puro.

\subsection{Bellman-Ford}

O algoritmo Bellman-Ford \cite{bellman} resolve caminhos m\'{\i}nimos em grafos com arestas de peso arbitr\'ario. A implementa\c{c}\~ao manual executa $V-1$ rodadas de relaxa\c{c}\~ao sobre todas as arestas:

\begin{equation}
\text{se } d[u] + w(u, v) < d[v] \text{ entonces } d[v] \leftarrow d[u] + w(u, v)
\label{eq:relaxacao}
\end{equation}

A itera\c{c}\~ao adicional $V$ detecta ciclos de peso negativo. A complexidade \'e $O(V \cdot E)$.

\subsection{Tr\^ansito Sint\'etico}

O tr\^ansito \'e modelado por fatores de congestionamento aplicados ao tempo base de cada aresta, conforme a TABELA~II.

\begin{table}[htbp]
\centering
\caption{Fatores de Congestionamento por Tipo de Via}
\label{tab:transito}
\begin{tabular}{@{}ll@{}}
\toprule
\textbf{Categoria de via} & \textbf{Fator} \\ \midrule
Arterial (motorway, trunk, primary) & 2,0 \\
Coletora (secondary, tertiary) & 1,5 \\
Residencial (residential, living\_street) & 1,1 \\ \bottomrule
\end{tabular}
\end{table}

O tempo ajustado \'e dado por $t_{\text{traffic}} = t_{\text{base}} \times f_{\text{cat}}$, onde $f_{\text{cat}}$ \'e o fator correspondente \`a categoria da via.

%% ============================================================
%% IV. EXPERIMENTOS
%% ============================================================
\section{Experimentos}

\subsection{Configura\c{c}\~ao Experimental}

A malha vi\'aria foi obtida para [CIDADE] com raio de [RAIO] m a partir do ponto central. As estat\'{\i}sticas do grafo s\~ao: [DADOS] n\'os e [DADOS] arestas para a rede \texttt{drive}; [DADOS] n\'os e [DADOS] arestas para a rede \texttt{walk}. O n\'o de origem $A$ foi mapeado para [COORD\_A] e o destino $B$ para [COORD\_B] utilizando o m\'etodo \texttt{nearest\_nodes} do OSMnx. A dist\^ancia m\'axima de caminhada foi fixada em $X = [X]$ m.

\subsection{Procedimento}

O procedimento experimental consta das seguintes etapas:

\begin{enumerate}
    \item Download e processamento da malha vi\'aria (adi\c{c}\~ao de velocidades e tempos de percurso) via OSMnx;
    \item Mapeamento de $A$ e $B$ para os n\'os mais pr\'oximos nas redes \texttt{drive} e \texttt{walk};
    \item Identifica\c{c}\~ao de candidatos de embarque $P$ por meio de Dijkstra no grafo \texttt{walk}, com filtro $d(A, P) \le X$;
    \item Execu\c{c}\~ao dos quatro algoritmos para cada candidato $P$, com e sem tr\^ansito sint\'etico;
    \item Medi\c{c}\~ao de tempo de execu\c{c}\~ao, n\'os/arestas expandidos, dist\^ancia total e tempo total.
\end{enumerate}

%% ============================================================
%% V. RESULTADOS
%% ============================================================
\section{Resultados}

\subsection{Desempenho Computacional}

A TABELA~III resume o desempenho dos algoritmos para o par origem-destino $(A, B)$.

\begin{table}[htbp]
\centering
\caption{Desempenho dos Algoritmos (Par $A \to B$)}
\label{tab:desempenho}
\begin{tabular}{@{}lrrr@{}}
\toprule
\textbf{Algoritmo} & \textbf{Tempo (ms)} & \textbf{Arestas relax.} & \textbf{Dist. (s)} \\ \midrule
Dijkstra simples & [DADOS] & [DADOS] & [DADOS] \\
Dijkstra heap & [DADOS] & [DADOS] & [DADOS] \\
A* geogr\'afico & [DADOS] & [DADOS] & [DADOS] \\
Bellman-Ford & [DADOS] & [DADOS] & [DADOS] \\ \bottomrule
\end{tabular}
\end{table}

\subsection{Compara\c{c}\~ao de Rotas}

Os resultados de rota s\~ao apresentados na TABELA~IV.

\begin{table}[htbp]
\centering
\caption{Compara\c{c}\~ao de Rotas}
\label{tab:rotas}
\begin{tabular}{@{}ll@{}}
\toprule
\textbf{M\'etrica} & \textbf{Valor} \\ \midrule
Menor dist\^ancia & [DADOS] \\
Rota mais r\'apida (s/ tr\^ansito) & [DADOS] \\
Rota mais r\'apida (c/ tr\^ansito) & [DADOS] \\
Sem caminhada ($A \to B$ direto) & [DADOS] \\ \bottomrule
\end{tabular}
\end{table}

\subsection{Ganho com Caminhada}

O candidato \'otimo $P^*$ identificado foi [DADOS], resultando em um ganho de tempo de [DADOS] segundos em rela\c{c}\~ao ao embarque imediato em $A$.

%% ============================================================
%% VI. DISCUSSÃO CRÍTICA
%% ============================================================
\section{Discuss\~ao Cr\'{\i}tica}

A modelagem como grafo dirigido com pesos em arestas (Se\c{c}\~ao~II) permite representar adequadamente a assimetria das ruas de m\~ao \u\'nica e a heterogeneidade dos tempos de percurso. Os tres pesos utilizados (dist\^ancia, tempo e tempo com congestionamento) capturam dimens\~oes complementares do custo de viagem.

O tr\^ansito sint\'etico alterou significativamente as rotas: [DADOS]. A caminhada at\'e um ponto de embarque alternativo melhorou a solu\c{c}\~ao em [DADOS] dos casos analisados, embora tenha se mostrado contraproducente quando [DADOS].

A menor dist\^ancia nem sempre corresponde \`a rota mais r\'apida [DADOS], evidenciando a import\^ancia de modelar m\'ultiplas fun\c{c}\~oes de custo. O algoritmo A* expandiu [DADOS]\% menos n\'os que o Dijkstra simples, consistente com a propriedade admiss\'{\i}vel da heur\'{\i}stica euclidiana. O Dijkstra com heap apresentou ganho de [DADOS]$\times$ sobre a implementa\c{c}\~ao simples, justificando o uso da fila de prioridade.

O Bellman-Ford n\~ao apresentou ganho de desempenho neste cen\'ario (pesos estritamente positivos), por\'em sua capacidade de detec\c{c}\~ao de ciclos negativos o torna mais robusto para aplica\c{c}\~oes com tr\^ansito din\^amico.

As principais limita\c{c}\~oes desta modelagem s\~ao: (i) tr\^ansito baseado em fatores fixos por tipo de via, sem dados em tempo real; (ii) aus\^encia de eventos sazonais e hor\'arios de pico; (iii) heur\'{\i}stica geogr\'afica simplificada que n\~ao considera a geometria real das ruas.

%% ============================================================
%% VII. CONCLUSÃO
%% ============================================================
\section{Conclus\~ao}

Este trabalho apresentou a modelagem de uma malha vi\'aria real como grafo dirigido e a compara\c{c}\~ao de quatro algoritmos de caminho m\'{\i}nimo para o problema de sele\c{c}\~ao de pontos de embarque. Os principais resultados indicam que [DADOS, como algoritmo mais eficiente e m\'etricas de desempenho].

O tr\^ansito sint\'etico alterou a escolha do candidato $P$ em [DADOS] dos casos, e a caminhada at\'e um ponto alternativo revelou-se vantajosa em [DADOS] das situa\c{c}\~oes testadas, demonstrando o trade-off entre dist\^ancia percorrida a p\'e e tempo total de viagem.

Como trabalhos futuros, destaca-se a integra\c{c}\~ao com dados de tr\^ansito em tempo real, a extens\~ao para m\'ultiplas fun\c{c}\~oes de custo (consumo de combust\'ivel, emiss\~oes) e o desenvolvimento de uma API de mobilidade para aplica\c{c}\~oes pr\'aticas.

%% ============================================================
%% REFERÊNCIAS
%% ============================================================
\bibliographystyle{IEEEtran}

\begin{thebibliography}{6}

\bibitem{boeing}
G. Boeing, ``Urban spatial order: Street network orientation, configuration, and entropy,'' \emph{Applied Network Science}, vol.~4, no.~1, pp.~1--25, 2019.

\bibitem{dijkstra}
E.~W. Dijkstra, ``A note on two problems in connexion with graphs,'' \emph{Numerische Mathematik}, vol.~1, no.~1, pp.~269--271, 1959.

\bibitem{hart}
P.~E. Hart, N.~J. Nilsson, and B. Raphael, ``A formal basis for the heuristic determination of minimum cost paths,'' \emph{IEEE Transactions on Systems Science and Cybernetics}, vol.~SSC-4, no.~2, pp.~100--107, 1968.

\bibitem{bellman}
R. Bellman, ``On a routing problem,'' \emph{Quarterly of Applied Mathematics}, vol.~16, no.~1, pp.~87--90, 1958.

\bibitem{networkx}
A.~A. Hagberg, D.~A. Schult, and P.~J. Swart, ``Exploring network structure, dynamics, and function using NetworkX,'' in \emph{Proc. 7th Python in Science Conference}, 2008, pp.~11--15.

\bibitem{osmnx}
G. Boeing, ``OSMnx: New methods for acquiring, constructing, analyzing, and visualizing complex street networks,'' \emph{Computers, Environment and Urban Systems}, vol.~65, pp.~126--139, 2017.

\end{thebibliography}

\end{document}
